\NormalFont\ShortTitle{2 Timotei}

{\MT A doua scrisoare a lui Pavel către Timotei


\par }\ChapOne{1}{\PP \VerseOne{1}Pavel, apostol al lui Isus Hristos, prin voia lui Dumnezeu, după făgăduința vieții care este în Hristos Isus,

\VS{2}către Timotei, copilul meu preaiubit: Har, îndurare și pace, de la Dumnezeu Tatăl și de la Hristos Isus, Domnul nostru.


\par }{\PP \VS{3}Mulțumesc lui Dumnezeu, căruia îi slujesc cu conștiință curată, ca și strămoșii mei. Cât de neîncetat îmi amintesc de voi în cererile mele, zi și noapte

\VS{4}tânjind să vă văd, amintindu-mi de lacrimile voastre, ca să mă umplu de bucurie;

\VS{5}după ce mi-am adus aminte de credința sinceră care este în voi, care a trăit mai întâi în bunica voastră Lois și în mama voastră Eunice și, sunt convins, și în voi.


\par }{\PP \VS{6}De aceea vă reamintesc să treziți darul lui Dumnezeu, care este în voi, prin punerea mâinilor mele.

\VS{7}Căci Dumnezeu nu ne-a dat un duh de frică, ci de putere, de dragoste și de stăpânire de sine.

\VS{8}De aceea, să nu vă fie rușine de mărturia Domnului nostru, nici de mine, prizonierul lui, ci să îndurați greutățile pentru Buna Vestire, potrivit cu puterea lui Dumnezeu,

\VS{9}care ne-a salvat și ne-a chemat cu o chemare sfântă, nu după faptele noastre, ci după propriul său scop și har, care ne-a fost dat în Isus Hristos înainte de veșnicie,

\VS{10}dar care a fost descoperit acum prin apariția Mântuitorului nostru, Isus Hristos, care a desființat moartea și a adus la lumină viața și nemurirea prin Buna Vestire.

\VS{11}Pentru aceasta am fost desemnat ca predicator, apostol și învățător al neamurilor.

\VS{12}De aceea și eu sufăr aceste lucruri.

\par }{\PP Dar nu mă rușinez, pentru că știu pe Cel în care am crezut și sunt încredințat că El poate păzi ceea ce i-am încredințat pentru ziua aceea.


\par }{\PP \VS{13}Țineți modelul cuvintelor sănătoase pe care le-ați auzit de la mine, în credința și dragostea care este în Hristos Isus.

\VS{14}Acel lucru bun care v-a fost încredințat, păziți-l prin Duhul Sfânt care locuiește în noi.


\par }{\PP \VS{15}Știți că toți cei din Asia s-au depărtat de mine, printre care Fgel și Hermogene.

\VS{16}Domnul să dea îndurare casei lui Onisifor, pentru că de multe ori m-a răcorit și nu i-a fost rușine de lanțul meu,

\VS{17}ci, când era la Roma, m-a căutat cu sârguință și m-a găsit

\VS{18}(Domnul să-i dea să găsească îndurarea Domnului în ziua aceea); și în câte lucruri a slujit la Efes, știți foarte bine.



\par }\Chap{2}{\PP \VerseOne{1}Tu, deci, copilul meu, întărește-te în harul care este în Hristos Isus.

\VS{2}Lucrurile pe care le-ai auzit de la mine printre mulți martori, încredințează-le unor oameni credincioși, care vor putea să învețe și pe alții.

\VS{3}De aceea, trebuie să înduri greutăți ca un bun soldat al lui Hristos Isus.

\VS{4}Nici un soldat aflat la datorie nu se încurcă în treburile vieții, pentru a fi pe placul celui care l-a înrolat ca soldat.

\VS{5}De asemenea, dacă cineva concurează în atletism, nu este încoronat decât dacă a concurat conform regulilor.

\VS{6}Fermierul care muncește trebuie să fie primul care primește o parte din recoltă.

\VS{7}Luați în considerare ceea ce vă spun și Domnul să vă dea pricepere în toate lucrurile.


\par }{\PP \VS{8}Adu-ți aminte de Iisus Hristos, înviat din morți, din neamul lui David, după Buna Vestire a mea,

\VS{9}în care am suferit până la lanțuri ca un criminal. Dar cuvântul lui Dumnezeu nu este înlănțuit.

\VS{10}De aceea îndur toate lucrurile de dragul celor aleși, pentru ca și ei să dobândească mântuirea care este în Cristos Isus cu glorie veșnică.

\VS{11}Acest cuvânt este demn de încredere:

\par }{\Q “Căci dacă am murit împreună cu el,

\par }{\QB vom trăi și noi cu el.


\par }{\Q \VS{12}Dacă rezistăm,

\par }{\QB vom domni și noi împreună cu el.

\par }{\Q Dacă îl negăm,

\par }{\QB ne va refuza și el.


\par }{\Q \VS{13}Dacă suntem necredincioși,

\par }{\QB el rămâne credincios;

\par }{\QB pentru că nu se poate nega pe sine.”


\par }{\PP \VS{14}Amintește-le aceste lucruri, cerându-le, înaintea Domnului, să nu se certe în zadar, ca să nu se certe în zadar, ca să strice pe cei ce ascultă.


\par }{\PP \VS{15}Străduiește-te să te prezinți pe tine însuți aprobat de Dumnezeu, ca un lucrător care nu are de ce să se rușineze, care să mânuiască bine Cuvântul adevărului.

\VS{16}Dar ferește-te de vorbăria goală, pentru că aceasta va merge mai departe în impietate,

\VS{17}iar cuvintele acelea vor mistui ca o cangrenă, dintre care se numără Himeneu și Filet:

\VS{18}oameni care au greșit în privința adevărului, spunând că învierea a trecut deja și răsturnând credința unora.

\VS{19}Cu toate acestea, temelia fermă a lui Dumnezeu rămâne în picioare, având această pecete: “Domnul cunoaște pe cei ce sunt ai lui” și: “Oricine numește numele Domnului să se depărteze de nedreptate”.


\par }{\PP \VS{20}Într-o casă mare nu sunt numai vase de aur și de argint, ci și de lemn și de lut. Unele sunt pentru cinste și altele pentru dezonoare.

\VS{21}Așadar, dacă cineva se curăță de acestea, va fi un vas de cinste, sfințit și potrivit pentru uzul stăpânului, pregătit pentru orice lucrare bună.


\par }{\PP \VS{22}Fugiți de poftele tinereții, și căutați dreptatea, credința, dragostea și pacea cu cei ce cheamă pe Domnul dintr-o inimă curată.

\VS{23}Dar refuzați întrebările nebunești și ignorante, știind că ele generează certuri.

\VS{24}Slujitorul Domnului nu trebuie să se certe, ci să fie blând cu toți, capabil să învețe, răbdător,

\VS{25}corectând cu blândețe pe cei care i se opun. Poate că Dumnezeu le va da pocăința care să ducă la deplina cunoaștere a adevărului

\VS{26}și se vor putea recupera din capcana diavolului, după ce au fost luați prizonieri de el pentru a-i face voia.



\par }\Chap{3}{\PP \VerseOne{1}Dar să știți că în zilele din urmă vor veni vremuri grele.

\VS{2}Căci oamenii vor fi iubitori de sine, iubitori de bani, lăudăroși, aroganți, hulitori, neascultători de părinți, nerecunoscători, necredincioși, nesfânți,

\VS{3}lipsiți de afecțiune firească, neiertători, calomniatori, fără stăpânire de sine, înverșunați, nu iubitori de bine,

\VS{4}trădători, încăpățânați, îngâmfați, încrezuți, iubitori de plăceri mai degrabă decât iubitori de Dumnezeu,

\VS{5}având o formă de evlavie, dar negând puterea ei. Îndepărtați-vă și de aceștia.

\VS{6}Căci unii dintre aceștia sunt oameni care se strecoară în case și iau prizoniere femei creduli, încărcate de păcate, duse de felurite pofte,

\VS{7}învățând mereu și neputând niciodată să ajungă la cunoașterea adevărului.

\VS{8}Așa cum Jannes și Jambres s-au împotrivit lui Moise, tot așa și aceștia se împotrivesc adevărului, oameni corupți la minte, care în privința credinței sunt respinși.

\VS{9}Dar ei nu vor merge mai departe. Căci nebunia lor va fi evidentă pentru toți oamenii, așa cum a ajuns să fie și a lor.


\par }{\PP \VS{10}Dar voi ați urmărit învățătura, purtarea, hotărârea, credința, răbdarea, dragostea, statornicia,

\VS{11}persecuțiile și suferințele mele, care mi s-au întâmplat în Antiohia, în Iconiu și în Listra. Am îndurat acele persecuții. Domnul m-a izbăvit din toate acestea.

\VS{12}Da, și toți cei care doresc să trăiască cu evlavie în Hristos Isus vor suferi persecuții.

\VS{13}Dar oamenii răi și impostorii se vor înrăutăți din ce în ce mai mult, înșelând și fiind înșelați.

\VS{14}Dar voi rămâneți în lucrurile pe care le-ați învățat și de care ați fost încredințați, știind de la cine le-ați învățat.

\VS{15}Încă din copilărie ați cunoscut Sfintele Scripturi, care pot să vă facă înțelepți pentru mântuire, prin credința care este în Hristos Isus.

\VS{16}Orice Scriptură este insuflată de Dumnezeu și folositoare pentru învățătură, pentru mustrare, pentru îndreptare și pentru instruire în neprihănire,

\VS{17}pentru ca fiecare persoană care aparține lui Dumnezeu să fie completă, echipată temeinic pentru orice lucrare bună.



\par }\Chap{4}{\PP \VerseOne{1}Vă poruncesc, așadar, înaintea lui Dumnezeu și a Domnului Isus Hristos, care va judeca pe cei vii și pe cei morți, la arătarea Lui și în Împărăția Lui:

\VS{2}propovăduiți Cuvântul; fiți grabnici la vreme și la nesfârșit; mustrați, mustrați și îndemnați cu toată răbdarea și cu toată învățătura.

\VS{3}Căci va veni vremea când nu vor asculta învățătura sănătoasă, ci, având mâncărimi de urechi, își vor îngrămădi învățători după poftele lor,

\VS{4}își vor întoarce urechile de la adevăr și se vor abate spre fabule.

\VS{5}Dar tu să fii cumpătat în toate, să suporți greutăți, să faci lucrarea de evanghelist și să-ți împlinești slujba.


\par }{\PP \VS{6}Căci deja sunt oferit și a venit vremea plecării mele.

\VS{7}Am luptat lupta cea bună. Am terminat cursul. Am păstrat credința.

\VS{8}De acum încolo, mi s-a strâns cununa dreptății, pe care Domnul, judecătorul cel drept, mi-o va da în ziua aceea; și nu numai mie, ci și tuturor celor care au iubit arătarea Lui.


\par }{\PP \VS{9}Să vă sârguiți să veniți curând la mine,

\VS{10}căci Demas m-a părăsit, iubind lumea de acum, și s-a dus la Tesalonic, Crescens în Galatia și Tit în Dalmația.

\VS{11}Numai Luca este cu mine. Luați-l pe Marcu și aduceți-l cu voi, căci îmi este de folos pentru slujbă.

\VS{12}Am trimis însă pe Tihicus la Efes.

\VS{13}Adu mantia pe care am lăsat-o la Troa la Carpus când vii — și cărțile, mai ales pergamentele.

\VS{14}Alexandru, coptarul, mi-a făcut mult rău. Domnul îi va răsplăti după faptele lui.

\VS{15}Ferește-te de el, pentru că s-a opus mult cuvintelor noastre.


\par }{\PP \VS{16}La prima mea apărare, nimeni nu a venit să mă ajute, ci toți m-au părăsit. Să nu li se reproșeze acest lucru.

\VS{17}Dar Domnul a stat lângă mine și m-a întărit, pentru ca prin mine să se proclame pe deplin mesajul și să audă toate neamurile. Astfel am fost izbăvit din gura leului.

\VS{18}Și Domnul mă va izbăvi de orice lucrare rea și mă va păstra pentru Împărăția Sa cerească. A lui să fie gloria în vecii vecilor. Amin.


\par }{\PP \VS{19}Salutați pe Prisca și pe Acuila și casa lui Onisifor.

\VS{20}Erast a rămas la Corint, dar pe Trofim l-am lăsat bolnav la Milet.

\VS{21}Fiți sârguincioși să veniți înainte de iarnă. Eubulus vă salută, ca și Pudens, Linus, Claudia și toți frații.


\par }{\PP \VS{22}Domnul Isus Hristos să fie cu duhul tău. Harul să fie cu tine. Amin.


\par }